% META: {"id":"TCOMPL-INEG-EV-B","chapitre":"TCOMPL-INEGALITES","type_objet":"evaluation","version":"B","duree_min":45,"bareme_total":20,"capacites_codes":["C1","C3","C5"],"competences":["calculer","raisonner"],"mode_creation":"ex_nihilo","sources_inspiration":[],"status":"approved","origin":"generated","status_history":["generated","audited","accepted","approved"],"release_acceptance":"RELEASE_OWNER_BATCH_ACCEPTANCE","acceptance_closure_digest":"sha256:9b3ccf9a81c5520fbb7e03b7b2d3e7bf2057a02834b6d908bf3be5f49f3b553f","acceptance_receipt_digest":"sha256:4fad86f0282e0fa4e0419cca1ad6379ae7af5a280c04a4a90880f90a1c044242"}
% BEGIN-VERIFY
% from sympy import symbols, diff, integrate, Rational
% x = symbols('x')
% L1 = (x**2 + 3*x) / 4
% assert L1.subs(x, 0) == 0 and L1.subs(x, 1) == 1
% I1 = integrate(L1, (x, 0, 1))
% G1 = 1 - 2*I1
% assert G1 == Rational(1, 12)
% L2 = x**2
% I2 = integrate(L2, (x, 0, 1))
% G2 = 1 - 2*I2
% assert G2 == Rational(1, 3)
% assert G2 > G1
% END-VERIFY

%---------------------------------------------------------------
% Duree : 45 minutes — Bareme : 20 points
%---------------------------------------------------------------

\begin{evaluation}{TCOMPL-INEG-EV-B}{45}

\begin{exercice}{TCOMPL-INEG-EV-B-EX1}{2}{40}
\textbf{Exercice unique} (20 points) --- \textit{Capacités C1, C3, C5}

\medskip

Deux pays A et B sont modélisés par les courbes de Lorenz $L_A(x) = \dfrac{x^2+3x}{4}$ et $L_B(x) = x^2$ sur $[0,1]$.
\begin{enumerate}
  \item (4 pts) Vérifier que $L_A(0)=0$ et $L_A(1)=1$.
  \item (6 pts) Calculer l'indice de Gini $G_A$ associé au pays A.
  \item (6 pts) Calculer l'indice de Gini $G_B$ associé au pays B (on rappelle $\displaystyle\int_0^1 x^2\,\mathrm{d}x=\dfrac{1}{3}$).
  \item (4 pts) Comparer $G_A$ et $G_B$ : quel pays a la répartition la plus inégalitaire ?
\end{enumerate}
\end{exercice}

\end{evaluation}
